\documentclass[a4paper,12pt]{article}
\usepackage[utf8]{inputenc}
\usepackage[english,russian]{babel}
\usepackage[T2A]{fontenc}

\usepackage[
  a4paper, mag=1000, includefoot,
  left=1.1cm, right=1.1cm, top=1.2cm, bottom=1.2cm, headsep=0.8cm, footskip=0.8cm
]{geometry}

\usepackage{amsmath}
\usepackage{amssymb}
\usepackage{times}
\usepackage{mathptmx}

\newcommand{\br}[1]{\left(#1\right)}

\IfFileExists{pscyr.sty}
{
  \usepackage{pscyr}
  \def\rmdefault{ftm}
  \def\sfdefault{ftx}
  \def\ttdefault{fer}
  \DeclareMathAlphabet{\mathbf}{OT1}{ftm}{bx}{it} % bx/it or bx/m
}

\mathsurround=0.1em
\clubpenalty=1000%
\widowpenalty=1000%
\brokenpenalty=2000%
\frenchspacing%
\tolerance=2500%
\hbadness=1500%
\vbadness=1500%
\doublehyphendemerits=50000%
\finalhyphendemerits=25000%
\adjdemerits=50000%


\begin{document}

\author{Коренев Иван}
\title{Метод Гаусса решения линейной системы с выбором главного элемента по всей матрице.(Многопоточная версия)}
\date{\today}
\maketitle

\section{Постановка задачи}
Решается система $Ax = y$ с вещественной квадратной матрицей коээфициентов размера $n$, где $x, y \in \mathbb{R}^n$ - столбцы неизвестных и свободных членов, соответственно, с помощью блочного многопоточного метода Гаусса с выбором главного элемента по всей матрице. Всего запусается $p$ потоков.

\section{Описание хранения данных}
Матрица хранится в одномерном массиве $A$ поблочно, сами блоки хранятся построчно. Для массива $A$ выделяется ровно $n^2$ элементов. Вектора правой части и неизвестных хранятся в одномерных массивах длины $n$. Координаты элемента матрицы $A$ в самой матрице называются глобальными, в блоке - локальными.

Поток с номером $t$, $0 \leq t < p$, работает с блочными строками, у которых номер по модулю $p$ равен $t$. Распределение данных приведено ниже. Объёмы массивов указаны в числе элементов. Собственные данные потока использует только этот поток. 

\textbf{Собственные данные каждого потока $t$}
\begin{itemize}
    \item Строка с главным элементом: $mn$.
    \item Обратная матрица для кандидата на главный элемент: $m^2$.
    \item Проверяемый блок: $m^2$.
    \item Лучший кандидат на обратный блок: $m^2$.
    \item Перестановки при обращении блока: $m$.
    \item Буфер для умножения матрицы на вектор: $m$.
    \item Индексы, норма лучшей матрицы, координаты кандидата: отдельные переменные.
\end{itemize}

Начальное заполнение общих массивов выполняется параллельно. Сначала поток $t$ инициализирует блочные строки матрицы $A$ с номерами $i \bmod p = t$ и обнуляет соответствующие блоки $x_i$. Элементы массива перестановок с номерами $0 \leq i < k$, $i \bmod p = t$, тот же поток инициализирует номерами $i$. После заполнения матрицы $A$ все потоки синхронизируются. Затем поток $t$ вычисляет блоки правой части $y_i$ по своим заполненным строкам матрицы $A$. После вычисления правой части все потоки повторно синхронизируются и переходят к решению системы. Параметры системы задаёт поток $0$ до начала совместных вычислений.

На шаге $\alpha$ обозначим через $t_\alpha = \alpha \bmod p$ владельца главной блочной строки, а через $t_*$ поток, чей кандидат выбран главным блоком. Перестановку блочных строк выполняет поток $0$, а перестановку блочных столбцов каждый поток выполняет в своих строках. Ниже указано, какие потоки записывают и используют общие данные.

\textbf{Общие данные всех потоков}
\begin{itemize}
    \item Матрица $A$: записывают и используют все потоки.
    \item Правая часть $y$: записывают и используют все потоки.
    \item Решение $x$: записывают и используют все потоки.
    \item Перестановки блочных столбцов: при инициализации записывают все потоки. При решении записывает и использует поток $0$.
    \item Обратная матрица главного блока: записывает поток $t_*$, использует поток $t_\alpha$.
    \item Параметры системы $n$, $m$, $p$: записывает поток $0$, используют все потоки.
\end{itemize}

Чтение общих данных после их изменения другим потоком выполняется после синхронизации. На время перестановки блочных строк потоком $0$ остальные потоки ожидают. Перед восстановлением исходного порядка неизвестных все потоки завершают обратный ход.

Все индексы и номер шага начинаются с $0$.
\subsection{Разбиение на блоки}
Пусть $n = m k + l, k = \lfloor \frac{n}{m} \rfloor, 0 \leq l < m$. При $l = 0$ блоки с номером $k$ отсутствуют. При $l > 0$ разбиение имеет вид:
\[
    A =
    \begin{pmatrix}
        A_{0, 0}^{m \times m} & \cdots & A_{0, k - 1}^{m \times m} & A_{0, k}^{m \times l}     \\
        \vdots & \ddots & \vdots & \vdots                    \\
        A_{k - 1, 0}^{m \times m} & \cdots & A_{k - 1, k - 1}^{m \times m} & A_{k - 1, k}^{m \times l} \\
        A_{k, 0}^{l \times m} & \cdots & A_{k, k - 1}^{l \times m}. & A_{k, k}^{l \times l}
    \end{pmatrix}.
\]

При $k = 0$ имеется только блок $A_{0, 0}^{l \times l}$.

Размер блока в одном направлении равен $\operatorname{size}(i) = \min(m, n - i m)$. Связь локальных и глобальных координат имеет вид:
\[
    \left(A_{i, j}^{\operatorname{size}(i) \times \operatorname{size}(j)} \right)_{p, q} = a_{i m + p, j m + q},
\]
\[
    0 \leq i, j < \left\lceil \frac{n}{m} \right\rceil, 0 \leq p < \operatorname{size}(i), 0 \leq q < \operatorname{size}(j).
\]

Элемент $(p, q)$ блока с номерами $(i, j)$ находится в массиве $A$ по индексу
\[
    m (i n + j \operatorname{size}(i)) + p \operatorname{size}(j) + q,
\]
с теми же множеством индексов.

Номер шага указывается в скобках справа сверху.

\section{Описание функций}
В функциях $i, j$ - номера блока, $W$ - блок с этими координатам, размеры блока равны $wr = \operatorname{size}(i)$, $wc = \operatorname{size}(j)$. Элемент с локальными координатами $(p, q)$ хранится в $W$ по индексу $p$, $wc + q$.

\begin{enumerate}
    \item \texttt{get\_block(i, j, dest)} копирует $wr \times wc$ элементов из массива $A$, начиная с индекса $i m n + j m wr$, в массив \texttt{dest}. Матрица $A$ не изменяется.
    \item \texttt{set\_block(i, j, src)} копирует $wr \times wc$ элементов из массива \texttt{src} в массив $A$, начиная с индекса $i m n + j m wr$. Остальные блоки не изменяются.
    \item \texttt{get\_block\_address(i, j)} возвращает указатель на первый элемент блока, то есть $A + i m n + j m wr$. Копирование блока не выполняется. Изменение элементов по этому указателю изменяет соответствующий блок матрицы $A$.
\end{enumerate}

\section{Формулы метода}
\subsection{Прямой ход метода Гаусса}

На шаге $\alpha = 0, \dots, k - 1$ поток $t$, $0 \leq t < p$, выбирает кандидата:
\[
    (i_t, j_t) \in \underset
    {\substack{\alpha \leq i,j < k \\ i \bmod p = t \\
    \det\br{\br{A_{i,j}^{m \times m}}^{(\alpha)}} \neq 0}}
    {\arg\min}
    \left\| \br{\br{A_{i,j}^{m \times m}}^{(\alpha)}}^{-1} \right\|.
\]
Здесь $\|B\| = \sum b_{ij}^2$. При равенстве значений выбираются наименьшие координаты. Обратная матрица кандидата сохраняется.
\[
    \nu_t =
    \begin{cases}
        \left\|\br{\br{A_{i_t,j_t}^{m \times m}}^{(\alpha)}}^{-1}\right\|,
            & \text{если кандидат найден}, \\
        -1, & \text{если кандидат не найден}.
    \end{cases}
\]

Общий выбор:
\[
    \nu_* =
    \begin{cases}
        \displaystyle\min_{\substack{0 \leq t < p \\ \nu_t > 0}}\nu_t,
            & \text{если найден хотя бы один кандидат}, \\
        -1, & \text{если кандидатов нет у всех потоков}.
    \end{cases}
\]
При $\nu_* > 0$:
\[
    t_* \in \{t \mid 0 \leq t < p, \nu_t = \nu_*\},
    (i_*, j_*) = (i_{t_*}, j_{t_*}).
\]
При равенстве значений выбираются наименьшие координаты. При $\nu_* = -1$ алгоритм завершается с отказом.

\textbf{Точка синхронизации после выбора главного блока.}

Поток $t_*$ записывает обратную матрицу кандидата в $L$. Поток $0$ переставляет блочные строки $\alpha$ и $i_*$ и соответствующие блоки $y$.

\textbf{Точка синхронизации после записи $L$ и перестановки строк.}

Поток $t$ переставляет блочные столбцы $\alpha$ и $j_*$ в строках $0 \leq i < \lceil n/m \rceil$, $i \bmod p = t$. Поток $0$ запоминает перестановку столбцов.

\textbf{Точка синхронизации после перестановки столбцов.}

После перестановок при $l > 0$:
\[
    \begin{pmatrix}
        E^{m \times m} & \br{A_{0, 1}^{m \times m}}^{(\alpha)} & \br{A_{0, 2}^{m \times m}}^{(\alpha)} & \dots  & \br{A_{0, \alpha}^{m \times m}}^{(\alpha)}          & \dots  & \br{A_{0, k - 1}^{m \times m}}^{(\alpha)}          & \br{A_{0, k}^{m \times l}}^{(\alpha)}          \\
        0              & E^{m \times m}                        & \br{A_{1, 2}^{m \times m}}^{(\alpha)} & \dots  & \br{A_{1, \alpha}^{m \times m}}^{(\alpha)}          & \dots  & \br{A_{1, k - 1}^{m \times m}}^{(\alpha)}          & \br{A_{1, k}^{m \times l}}^{(\alpha)}          \\
        \vdots         & \vdots                                & \vdots                                & \ddots & \vdots                                              & \vdots & \vdots                                             & \vdots                                         \\
        0              & \dots                                 & 0                                     & \dots  & \br{A_{\alpha, \alpha}^{m \times m}}^{(\alpha)}     & \dots  & \br{A_{\alpha, k - 1}^{m \times m}}^{(\alpha)}     & \br{A_{\alpha, k}^{m \times l}}^{(\alpha)}     \\
        0              & \dots                                 & 0                                     & \dots  & \br{A_{\alpha + 1, \alpha}^{m \times m}}^{(\alpha)} & \dots  & \br{A_{\alpha + 1, k - 1}^{m \times m}}^{(\alpha)} & \br{A_{\alpha + 1, k}^{m \times l}}^{(\alpha)} \\
        \vdots         & \vdots                                & \vdots                                & \ddots & \vdots                                              & \vdots & \vdots                                             & \vdots                                         \\
        0              & \dots                                 & 0                                     & \dots  & \br{A_{k - 1, \alpha}^{m \times m}}^{(\alpha)}      & \dots  & \br{A_{k - 1, k - 1}^{m \times m}}^{(\alpha)}      & \br{A_{k - 1, k}^{m \times l}}^{(\alpha)}      \\
        0              & \dots                                 & 0                                     & \dots  & \br{A_{k, \alpha}^{l \times m}}^{(\alpha)}          & \dots  & \br{A_{k, k - 1}^{l \times m}}^{(\alpha)}          & \br{A_{k, k}^{l \times l}}^{(\alpha)}
    \end{pmatrix}.
\]

При $l = 0$ последняя блочная строка и последний блочный столбец в этой записи отсутствуют.

Пусть $L = \br{\br{A_{\alpha, \alpha}^{m \times m}}^{(\alpha)}}^{-1}$. Нормировка потоком $t_\alpha = \alpha \bmod p$:
\begin{align*}
    \br{A_{\alpha, \alpha}^{m \times m}}^{(\alpha + 1)} & = E^{m \times m},
        t = t_\alpha, \\
    \br{A_{\alpha, j}^{m \times m}}^{(\alpha + 1)} & = L \br{A_{\alpha, j}^{m \times m}}^{(\alpha)},
        j = \alpha + 1, \dots, k - 1,t = t_\alpha, \\
    \br{y_{\alpha}^{m \times 1}}^{(\alpha + 1)} & = L \br{y_{\alpha}^{m \times 1}}^{(\alpha)},
        t = t_\alpha, \\
    \br{A_{\alpha, k}^{m \times l}}^{(\alpha + 1)} & = L \br{A_{\alpha, k}^{m \times l}}^{(\alpha)},
        l > 0,t = t_\alpha.
\end{align*}
\textbf{Точка синхронизации после нормировки главной строки.}

Каждый поток копирует главную строку в свой буфер.

Полные блочные строки, $i = \alpha + 1, \dots, k - 1$, $i \bmod p = t$:
\begin{align*}
    \br{A_{i, j}^{m \times m}}^{(\alpha + 1)} & = \br{A_{i, j}^{m \times m}}^{(\alpha)}
        - \br{A_{i, \alpha}^{m \times m}}^{(\alpha)} \br{A_{\alpha, j}^{m \times m}}^{(\alpha + 1)},
        j = \alpha + 1, \dots, k - 1, \\
    \br{A_{i, k}^{m \times l}}^{(\alpha + 1)} & = \br{A_{i, k}^{m \times l}}^{(\alpha)}
        - \br{A_{i, \alpha}^{m \times m}}^{(\alpha)} \br{A_{\alpha, k}^{m \times l}}^{(\alpha + 1)},
        l > 0, \\
    \br{y_i^{m \times 1}}^{(\alpha + 1)} & = \br{y_i^{m \times 1}}^{(\alpha)}
        - \br{A_{i, \alpha}^{m \times m}}^{(\alpha)} \br{y_\alpha^{m \times 1}}^{(\alpha + 1)}, \\
    \br{A_{i, \alpha}^{m \times m}}^{(\alpha + 1)} & = 0.
\end{align*}
Неполная блочная строка, $l > 0$, $t = k \bmod p$:
\begin{align*}
    \br{A_{k, j}^{l \times m}}^{(\alpha + 1)} & = \br{A_{k, j}^{l \times m}}^{(\alpha)}
        - \br{A_{k, \alpha}^{l \times m}}^{(\alpha)} \br{A_{\alpha, j}^{m \times m}}^{(\alpha + 1)},
        j = \alpha + 1, \dots, k - 1, \\
    \br{A_{k, k}^{l \times l}}^{(\alpha + 1)} & = \br{A_{k, k}^{l \times l}}^{(\alpha)}
        - \br{A_{k, \alpha}^{l \times m}}^{(\alpha)} \br{A_{\alpha, k}^{m \times l}}^{(\alpha + 1)}, \\
    \br{y_k^{l \times 1}}^{(\alpha + 1)} & = \br{y_k^{l \times 1}}^{(\alpha)}
        - \br{A_{k, \alpha}^{l \times m}}^{(\alpha)} \br{y_\alpha^{m \times 1}}^{(\alpha + 1)}, \\
    \br{A_{k, \alpha}^{l \times m}}^{(\alpha + 1)} & = 0.
\end{align*}
Блок $A_{i,\alpha}^{(\alpha)}$ зануляется после обновления строки и блока $y_i$.

\textbf{Точка синхронизации после обновления строк.}

После $k$ шагов при $l > 0$:
\[
    \begin{pmatrix}
        E^{m \times m} & \br{A_{0, 1}^{m \times m}}^{(k)}     & \br{A_{0, 2}^{m \times m}}^{(k)}     & \dots  & \br{A_{0, k - 1}^{m \times m}}^{(k)}     & \br{A_{0, k}^{m \times l}}^{(k)}        & \br{y_{0}^{m \times 1}}^{(k)}     \\
        0              & E^{m \times m}                       & \br{A_{1, 2}^{m \times m}}^{(k)}     & \dots  & \br{A_{1, k - 1}^{m \times m}}^{(k)}     & \br{A_{1, k}^{m \times l}}^{(k)}        & \br{y_{1}^{m \times 1}}^{(k)}     \\
        \vdots         & \vdots                               & \vdots                               & \ddots & \vdots                                   & \vdots                                  & \vdots                            \\
        0              & 0                                    &  0                                   & \dots  & E^{m \times m}                           & \br{A_{k - 1, k}^{m \times l}}^{\br{k}} & \br{y_{k - 1}^{m \times 1}}^{(k)} \\
        0              & 0                                    &  0                                   & \dots  & 0                                        & \br{A_{k, k}^{l \times l}}^{(k)}        & \br{y_{k}^{l \times 1}}^{(k)}
    \end{pmatrix}.
\]

При $l = 0$ прямой ход завершён.

При $l > 0$ поток $t = k \bmod p$ решает последнюю систему методом Гаусса с полным выбором:
\[
    \br{A_{k,k}^{l \times l}}^{(k)} \br{y_k^{l \times 1}}^{(k + 1)}
    = \br{y_k^{l \times 1}}^{(k)},
    l > 0,t = k \bmod p.
\]
\textbf{Точка синхронизации после решения последней системы.}

\subsection{Обратный ход метода Гаусса}
Закрепление строк за потоками сохраняется. Блочную строку $i$ обрабатывает поток $t = i \bmod p$, как в прямом ходе. Этот поток обновляет $y_i$ и записывает $x_i$.

Индексы прямого хода опускаем. Начальное состояние правой части после прямого хода и решения последней системы обозначаем $y^{(k)}$ при $l \neq 0$ и $y^{(k - 1)}$ при $l = 0$. Номер шага обратного хода убывает.

Находим последний блок решения:
\begin{align*}
    x_k^{l \times 1}         & = \br{y_k^{l \times 1}}^{(k)},
        l \neq 0, t = k \bmod p, \\
    x_{k - 1}^{m \times 1}   & = \br{y_{k - 1}^{m \times 1}}^{(k - 1)},
        l = 0, t = (k - 1) \bmod p.
\end{align*}
\textbf{Точка синхронизации после записи последнего блока решения.}

На шаге $\alpha = k, \dots, 1$ при $l \neq 0$ и $\alpha = k - 1, \dots, 1$ при $l = 0$ поток $t$ обновляет блоки правой части в строках $i = 0, \dots, \alpha - 1$, $i \bmod p = t$:
\begin{align*}
    \br{y_i^{m \times 1}}^{(\alpha - 1)} & = \br{y_i^{m \times 1}}^{(\alpha)}
        - A_{i,k}^{m \times l}x_k^{l \times 1},
        \alpha = k, l \neq 0, i \bmod p = t, \\
    \br{y_i^{m \times 1}}^{(\alpha - 1)} & = \br{y_i^{m \times 1}}^{(\alpha)}
        - A_{i,\alpha}^{m \times m}x_\alpha^{m \times 1},
        \alpha < k, i \bmod p = t.
\end{align*}
После обновления строки $\alpha - 1$ её владелец записывает следующий блок решения:
\[
    x_{\alpha - 1}^{m \times 1} = \br{y_{\alpha - 1}^{m \times 1}}^{(\alpha - 1)},
    t = (\alpha - 1) \bmod p.
\]
\textbf{Точка синхронизации после обновления строк и записи следующего блока решения.}

При единственном блоке шаги обновления не выполняются. Для восстановления исходного порядка неизвестных поток $0$ применяет к блокам вектора $x$ сохранённые перестановки столбцов в обратном порядке.

\section{Оценка числа точек синхронизации}
Пусть $n = mk + l$, $0 \leq l < m$, $1 \leq m \leq n$, $p \mid k$.

При $l = 0$:
\begin{align*}
    B_{\text{нач}} & = 2, \\
    B_{\text{пр}} & = 5k, \\
    B_{\text{обр}} & = 1 + (k - 1) = k.
\end{align*}
При $l > 0$ добавляются одна точка после решения последней системы и один шаг обратного хода. Итого:
\[
    B_{\max}(n,m,p) =
    \begin{cases}
        6k + 2, & l = 0, \\
        6k + 4, & l > 0.
    \end{cases}
\]

\section{Число арифметических операций}
Пусть $m \mid n$, $1 \leq m \leq n$, $k = n/m$, $p \mid k$, $q = k/p$. Кандидаты обратимы, выбранная обратная матрица сохраняется. Сравнения и перестановки не учитываем.

Через $S(n,m,p)$ обозначаем загрузку самого загруженного потока за прямой и обратный ход. Максимум достигается при $t = p - 1$. Этому потоку принадлежат строки $i = jp - 1$, $j = 1, \dots, q$. Далее считаем операции этого потока.

\subsection{В прямом ходе}
Обращение блока методом Гаусса с полным выбором и $m$ правыми частями:
\[
    I(m) = \sum_{r = 0}^{m - 1}\left[2r^2 + (4m + 1)r + m\right]
    = \frac{8}{3}m^3 - \frac{3}{2}m^2 - \frac{1}{6}m.
\]

Нормировка и исключение, $r_\alpha = n - (\alpha + 1)m$. На шаге $\alpha$ поток обновляет $\lceil(k - \alpha - 1)/p\rceil$ нижних строк и нормирует главную строку при $\alpha \bmod p = p - 1$:
\begin{align*}
    S_{\text{обн}}
    & = \sum_{j = 1}^{q}\left[m(r_{jp - 1} + 1)(2m - 1)
        + 2m^2\sum_{\alpha = 0}^{jp - 2}(r_\alpha + 1)\right] = \\
    & = \frac{\frac{2}{3}n^3 + \frac{p - 1}{2}n^2m
        - \frac{p(p + 3)}{6}nm^2}{p} + O(n^2).
\end{align*}
Число кандидатов и затраты на поиск:
\begin{align*}
    N & = \sum_{\alpha = 0}^{k - 1}(k - \alpha)
        \left\lceil\frac{k - \alpha}{p}\right\rceil = \\
    & = \frac{n^3}{3m^3p} + \frac{(p + 1)n^2}{4m^2p}
        + \frac{(3 - p)n}{12m}.
\end{align*}
\begin{align*}
    S_{\text{поиск}} & = N\left[I(m) + 2m^2 - 1\right] = \\
    & = \frac{1}{p}\Biggl[\frac{8}{9}n^3 + \frac{2(p + 1)}{3}n^2m
        + \frac{2p(3 - p)}{9}nm^2 + \\
    & + \frac{n^3}{6m} - \frac{n^3}{18m^2}
        - \frac{n^3}{3m^3}\Biggr] + O(n^2).
\end{align*}
Итого
\begin{align*}
    S_{\text{пр}} & = S_{\text{обн}} + S_{\text{поиск}} = \\
    & = \frac{1}{p}\Biggl[\frac{14}{9}n^3 + \frac{7p + 1}{6}n^2m
        - \frac{p(7p - 3)}{18}nm^2 + \\
    & + \frac{n^3}{6m} - \frac{n^3}{18m^2}
        - \frac{n^3}{3m^3}\Biggr] + O(n^2).
\end{align*}

\subsection{В обратном ходе}
\[
    S_{\text{обр}} = \sum_{j = 1}^{q}2m(n - jpm)
        = \frac{n^2}{p} - nm = O\br{\frac{n^2}{p}}.
\]

\subsection{Суммарное число}
\begin{align*}
    S(n,m,p) & = S_{\text{пр}} + S_{\text{обр}} = \\
    & = \frac{1}{p}\Biggl[\frac{14}{9}n^3 + \frac{7p + 1}{6}n^2m
        - \frac{p(7p - 3)}{18}nm^2 + \\
    & + \frac{n^3}{6m} - \frac{n^3}{18m^2}
        - \frac{n^3}{3m^3}\Biggr] + O(n^2).
\end{align*}

\section{Проверка формул}
\begin{align*}
    S(n,m,1) & = \frac{14}{9}n^3 + \frac{4}{3}n^2m - \frac{2}{9}nm^2 + \\
    & + \frac{n^3}{6m} - \frac{n^3}{18m^2} - \frac{n^3}{3m^3} + O(n^2), \\
    S(n,1,1) & = \frac{4}{3}n^3 + O(n^2), \\
    S(n,n,1) & = \frac{8}{3}n^3 + O(n^2).
\end{align*}
При $p = 1$ получаем формулу однопоточного метода. Для обычного метода Гаусса с полным выбором
\[
    S_{\text{Г}}(n) = \frac{2}{3}n^3 + O(n^2).
\]
При $m = 1$ старший коэффициент вдвое больше из-за вычисления обратного значения и его квадрата для каждого кандидата при поиске. При $m = n$ он вчетверо больше из-за вычисления всей обратной матрицы вместо решения одной системы.

\end{document}
